\documentclass[11pt]{article}
\usepackage[a4paper,margin=1in]{geometry}
\usepackage{amsmath,amssymb,amsthm}
\usepackage[hidelinks]{hyperref}

\newtheorem{theorem}{Theorem}
\newtheorem{lemma}{Lemma}
\newtheorem{corollary}{Corollary}
\newtheorem{remark}{Remark}

\newcommand{\K}{\mathbb K}
\newcommand{\ellone}{\ell_1}
\newcommand{\ellinf}{\ell_\infty}
\newcommand{\weakstar}{w^*}

\title{A Sharp One-Third Barrier for the \(M^*(n,n)\) Route in the Numerical-Radius \(M\)-Ideal Question}
\author{Automatic math-discovery packet}
\date{June 20, 2026}

\begin{document}
\maketitle

\section*{Source question}

The source paper is Manwook Han and Sun Kwang Kim,
\emph{Geometry of the space of compact operators endowed with the numerical radius norm},
arXiv:2502.10821.  For a Banach space \(X\) with positive numerical index \(n(X)\), the paper asks:

\begin{quote}
Let \(X\) be a Banach space with \(n(X)>0\). Is \(X\) Asplund whenever
\(\mathcal K_v(X)\) is an \(M\)-ideal in \(\mathcal L_v(X)\)?
\end{quote}

The relevant source mechanism is Proposition 3.7: if \(X\) is infinite-dimensional,
\(n(X)>0\), and \(\mathcal K_v(X)\) is an \(M\)-ideal in \(\mathcal L_v(X)\), then \(X\) has
property \(M^*(n(X),n(X))\).  The source then uses the Cabello--Nieto--Oja
theorem that the \(M(r,s)\)-inequality with \(r+s>1\) implies Asplundness, giving
the Asplund conclusion when \(n(X)>1/2\).  It also records that, under the same
\(n(X)>1/2\) hypothesis, a space containing \(\ell_1\) cannot satisfy the
numerical-radius \(M\)-ideal hypothesis.

\section*{Definitions}

For \(r,s\in(0,1]\), a Banach space \(X\) has property \(M^*(r,s)\) if, whenever
\(x^*,y^*\in X^*\), \(\|x^*\|\leq \|y^*\|\), and \((z_\beta^*)\) is a bounded
\(\weakstar\)-null net in \(X^*\), one has
\[
    \limsup_\beta \|r x^*+s z_\beta^*\|
    \leq
    \limsup_\beta \|y^*+z_\beta^*\|.
\]

\section*{Idea of the Proof}

The source proof loses the factor \(n(X)\) when it converts numerical radius
control back to operator-norm control.  The resulting property \(M^*(n,n)\)
has its own intrinsic threshold: \(M^*(r,r)\) is automatic for every Banach
space as soon as \(r\leq 1/3\).  This follows from the triangle inequality and
weak-star lower semicontinuity.

The constant is sharp.  In \(\ellinf=(\ellone)^*\), take the constant
functionals \(1\) and \(-1\), and the weak-star-null sequence \(2e_n\).  Then
\(\|1+2e_n\|_\infty=3\), while \(\|-1+2e_n\|_\infty=1\).  Thus \(M^*(r,r)\)
fails for every \(r>1/3\).  A standard weak-star lifting argument transfers
this obstruction from almost-isometric copies of \(\ellone\) inside an arbitrary
Banach space.

\section*{Main Result}

\begin{theorem}[Universal one-third bound]
\label{thm:universal}
Every Banach space \(X\) has property \(M^*(r,r)\) for every \(0<r\leq 1/3\).
The same proof gives property \(M(r,r)\) for every \(0<r\leq 1/3\).
\end{theorem}

\begin{proof}
We prove the \(M^*\) statement.  Let \(x^*,y^*\in X^*\) satisfy
\(\|x^*\|\leq\|y^*\|\), and let \((z_\beta^*)\) be a bounded \(\weakstar\)-null
net in \(X^*\).  Put
\[
    a=\|y^*\|,\qquad L=\limsup_\beta \|z_\beta^*\|.
\]
Since \(y^*+z_\beta^*\to y^*\) in the weak-star topology and the norm is
weak-star lower semicontinuous,
\[
    \limsup_\beta \|y^*+z_\beta^*\|\geq a.
\]
The triangle inequality also gives
\[
    \limsup_\beta \|y^*+z_\beta^*\|\geq L-a.
\]
Hence
\[
    \limsup_\beta \|y^*+z_\beta^*\|
    \geq
    \max\{a,L-a\}.
\]
For all \(a,L\geq0\),
\[
    \max\{a,L-a\}\geq \frac{a+L}{3}.
\]
Indeed, if \(L\leq 2a\), then \(a\geq(a+L)/3\); if \(L\geq2a\), then
\(L-a\geq(a+L)/3\).  On the other hand,
\[
    \limsup_\beta \|x^*+z_\beta^*\|
    \leq
    \|x^*\|+L
    \leq
    a+L.
\]
Therefore, for \(0<r\leq1/3\),
\[
    r\limsup_\beta \|x^*+z_\beta^*\|
    \leq
    \frac{a+L}{3}
    \leq
    \limsup_\beta \|y^*+z_\beta^*\|,
\]
which is exactly \(M^*(r,r)\).  The proof of \(M(r,r)\) is identical, replacing
\(\weakstar\)-null nets in \(X^*\) by weakly null nets in \(X\).
\end{proof}

\begin{theorem}[Sharpness at \(\ell_1\)]
\label{thm:sharp}
The space \(\ellone\) fails property \(M^*(r,r)\) for every \(r>1/3\).
More generally, if a Banach space \(X\) has property \(M^*(r,r)\) for some
\(r>1/3\), then \(X\) contains no isomorphic copy of \(\ellone\).
\end{theorem}

\begin{proof}
First take \(X=\ellone\), so \(X^*=\ellinf\) with its \(\sigma(\ellinf,\ellone)\)
topology.  Let
\[
    u=(1,1,1,\ldots),\qquad v=(-1,-1,-1,\ldots),\qquad z_n=2e_n.
\]
Then \(\|u\|_\infty=\|v\|_\infty=1\), and \(z_n\to0\) in
\(\sigma(\ellinf,\ellone)\).  Moreover
\[
    \|u+z_n\|_\infty=3,\qquad \|v+z_n\|_\infty=1
\]
for every \(n\).  Thus \(M^*(r,r)\) would force \(3r\leq1\).

We now explain the standard lifting argument for arbitrary \(X\).  Suppose that
\(X\) contains an isomorphic copy of \(\ellone\).  By James' distortion theorem
for \(\ellone\), for every \(D>1\) there is a subspace \(E\subset X\) and an
isomorphism \(U:\ellone\to E\) such that
\[
    \|a\|_1\leq \|Ua\|\leq D\|a\|_1\qquad (a\in\ellone).
\]
For \(b\in\ellinf\), let \(\phi_b\in E^*\) be given by
\(\phi_b(Ua)=\sum_k b_k a_k\).  Then
\[
    \|\phi_b\|_{E^*}\leq \|b\|_\infty,\qquad
    \|\phi_b\|_{E^*}\geq D^{-1}\|b\|_\infty.
\]
Let \(b^+=(1,1,\ldots)\), \(b^-=(-1,-1,\ldots)\), and
\(c_m=b^-+2e_m\).  The functionals \(\phi_{b^+}\) and \(\phi_{b^-}\) have the
same norm.  Choose norm-preserving Hahn--Banach extensions \(u,v\in X^*\) of
\(\phi_{b^+}\) and \(\phi_{b^-}\), respectively.

We use the following elementary finite-dimensional form of weak-star lifting.
For every finite-dimensional \(F\subset X\) and every \(\varepsilon>0\), for all
sufficiently large \(m\) there exists \(h_{m,F,\varepsilon}\in X^*\) extending
\(\phi_{c_m}\) such that
\[
    \|h_{m,F,\varepsilon}\|\leq 1+\varepsilon,\qquad
    |h_{m,F,\varepsilon}(x)-v(x)|\leq \varepsilon\|x\|
    \quad(x\in F).
\]
Indeed, \(\phi_{c_m}-\phi_{b^-}=2\phi_{e_m}\) converges pointwise to \(0\) on
\(E\), hence uniformly on the unit ball of the finite-dimensional space
\(E\cap F\).  The two prescriptions \(\phi_{c_m}\) on \(E\) and \(v|_F\) are
therefore compatible up to \(\varepsilon\) on \(E\cap F\); the usual
finite-dimensional Hahn--Banach perturbation lemma gives an extension on
\(E+F\), and Hahn--Banach extends it to \(X\).

Set \(z_{m,F,\varepsilon}=h_{m,F,\varepsilon}-v\), with the directed index
\((F,\varepsilon)\) ordered by inclusion of \(F\) and decreasing
\(\varepsilon\), and \(m\) chosen large enough as above.  Then
\((z_{m,F,\varepsilon})\) is a bounded \(\weakstar\)-null net in \(X^*\).
Furthermore,
\[
    \limsup \|v+z_{m,F,\varepsilon}\|
    =
    \limsup \|h_{m,F,\varepsilon}\|
    \leq 1,
\]
while the restriction of \(u+z_{m,F,\varepsilon}\) to \(E\) is
\(\phi_{b^+}-\phi_{b^-}+\phi_{c_m}=\phi_{b^+ + 2e_m}\), so
\[
    \limsup \|u+z_{m,F,\varepsilon}\|
    \geq
    D^{-1}\|b^+ + 2e_m\|_\infty
    =
    3D^{-1}.
\]
If \(X\) had property \(M^*(r,r)\), we would get \(3rD^{-1}\leq1\), i.e.
\(r\leq D/3\).  Since \(D>1\) is arbitrary, \(r\leq1/3\).  Hence no space with
property \(M^*(r,r)\), \(r>1/3\), can contain an isomorphic copy of \(\ellone\).
\end{proof}

\begin{corollary}[Improved \(\ell_1\)-exclusion for the source question]
\label{cor:source}
Let \(X\) be a Banach space with \(n(X)>1/3\).  If
\(\mathcal K_v(X)\) is an \(M\)-ideal in \(\mathcal L_v(X)\), then \(X\)
contains no isomorphic copy of \(\ellone\).
\end{corollary}

\begin{proof}
By Proposition 3.7 of Han--Kim, the numerical-radius \(M\)-ideal hypothesis
implies property \(M^*(n(X),n(X))\).  Since \(n(X)>1/3\), Theorem
\ref{thm:sharp} excludes every isomorphic copy of \(\ellone\).
\end{proof}

\begin{corollary}[Complex spaces]
If \(X\) is complex and \(\mathcal K_v(X)\) is an \(M\)-ideal in
\(\mathcal L_v(X)\), then \(X\) contains no isomorphic copy of \(\ellone\).
\end{corollary}

\begin{proof}
The standard Bonsall--Duncan lower bound gives \(n(X)\geq 1/e>1/3\) for every
complex Banach space.  Apply Corollary \ref{cor:source}.
\end{proof}

\section*{Limitations}

This packet does not prove that the source question has a positive answer.
The new conclusion rules out \(\ellone\)-containment under the stronger
threshold \(n(X)>1/3\), but Asplundness is stronger than not containing
\(\ellone\).  In particular, James-tree-type mechanisms show why
non-Asplund spaces without \(\ellone\) remain a genuine obstacle.

The packet also shows a hard limitation of the source paper's immediate route:
for \(n(X)\leq1/3\), property \(M^*(n(X),n(X))\) is automatic for every Banach
space, so it cannot by itself imply Asplundness.

\section*{Novelty and verification notes}

The cheap run indexes had no exact prior packet for arXiv:2502.10821 or this
numerical-radius \(M\)-ideal question.  A bounded web search for the exact
question and for phrases involving \(M^*(r,s)\), \(1/3\), and \(\ell_1\) found
the source paper but no later answer or matching constant-barrier result.

Human review should focus on the finite-dimensional lifting paragraph in
Theorem \ref{thm:sharp}; the remaining estimates are scalar norm calculations.
No numerical computation was used.

\begin{thebibliography}{9}

\bibitem{BD}
F. F. Bonsall and J. Duncan,
\emph{Numerical ranges of operators on normed spaces and of elements of normed algebras},
London Mathematical Society Lecture Note Series, vol. 2, Cambridge University Press, 1971.

\bibitem{CNO}
J. C. Cabello, E. Nieto, and E. Oja,
\emph{On ideals of compact operators satisfying the \(M(r,s)\)-inequality},
J. Math. Anal. Appl. 220 (1998), 334--348.

\bibitem{HanKim}
M. Han and S. K. Kim,
\emph{Geometry of the space of compact operators endowed with the numerical radius norm},
arXiv:2502.10821.

\bibitem{James}
R. C. James,
\emph{Uniformly non-square Banach spaces},
Ann. of Math. (2) 80 (1964), 542--550.

\end{thebibliography}

\end{document}
